
\documentclass{article}
\usepackage{colortbl}
\usepackage{makecell}
\usepackage{multirow}
\usepackage{supertabular}

\begin{document}

\newcounter{utterance}

\centering \large Interaction Transcript for game `function\_detective', experiment `numbers\_passive\_labeled\_pairs', episode 27 with mistral{-}large{-}3{-}azure.
\vspace{24pt}

{ \footnotesize  \setcounter{utterance}{1}
\setlength{\tabcolsep}{0pt}
\begin{supertabular}{c@{$\;$}|p{.15\linewidth}@{}p{.15\linewidth}p{.15\linewidth}p{.15\linewidth}p{.15\linewidth}p{.15\linewidth}}
    \# & \multicolumn{2}{c}{Player} && \multicolumn{2}{c}{Game Master} \\
    \hline

    \theutterance \stepcounter{utterance}  
    & & & \multicolumn{4}{p{0.6\linewidth}}{
        \cellcolor[rgb]{0.9,0.9,0.9}{
            \makecell[{{p{\linewidth}}}]{
                \texttt{\tiny{[P1$\langle$GM]}}
                \texttt{I have a hidden Python function for you to figure out.} \\
\\ 
\texttt{Your job is to infer what it does and then submit a correct Python implementation.} \\
\\ 
\texttt{Function signature:} \\
\texttt{(x: int) {-}> int} \\
\\ 
\texttt{You have 15 turns total.} \\
\\ 
\texttt{I reveal candidate input{-}output pairs one at a time and tell you whether each one is correct.} \\
\\ 
\texttt{You must make your best SOLVE attempt no later than the second{-}to{-}last turn.} \\
\texttt{Do not wait until the very end.} \\
\\ 
\texttt{You can use these actions:} \\
\texttt{{-} JSON format described at the end of the prompt} \\
\texttt{{-} SOLVE: ```python ... ```} \\
\\ 
\texttt{Do not add import statements in your guess.} \\
\texttt{Anything you might need is already available, so just write the function body in hypothesis and return the answer.} \\
\\ 
\texttt{Preloaded examples:} \\
\texttt{1. ((2, 2), True)} \\
\\ 
\texttt{Scoring:} \\
\texttt{{-} correctness is required for any non{-}zero score} \\
\texttt{{-} in interactive modes, solving with fewer information requests is better} \\
\\ 
\texttt{Input formatting reminders:} \\
\texttt{{-} strings must be quoted} \\
\texttt{{-} booleans must be True/False} \\
\texttt{{-} use commas between multiple values} \\
\texttt{{-} do not write any extra action tags in the same message} \\
\\ 
\texttt{Format example:} \\
\texttt{You write:} \\
\texttt{\{'current\_mode': <probing/confirming/don't know; choose whichever best describes your current state in the task>, 'input\_rationale': <brief description of why you chose this hypothesis>, 'hypothesis': <either the exact string "unknown" if you have no guess yet OR your current best guess as runnable Python matching the given signature, with a one{-}line docstring stating what it does. Never write placeholder code, if you cannot commit to a guess, use "unknown". Match the given signature.} \\
\texttt{e.g.} \\
\texttt{def f(x):} \\
\texttt{	"""Returns the absolute value of x."""} \\
\texttt{	return x if x >= 0 else {-}x} \\
\\ 
\texttt{Use "unknown" if you have no guess.>, 'confidence\_score': <score between 1{-}100. 100 means highest confidence in your hypothesis>\}} \\
\\ 
\texttt{Output only one valid json per turn.} \\
            }
        }
    }
     \\ \\

    \theutterance \stepcounter{utterance}  
    & \multicolumn{4}{p{0.6\linewidth}}{
        \cellcolor[rgb]{0.561,0.737,0.561}{
            \makecell[{{p{\linewidth}}}]{
                \texttt{\tiny{[P1$\rangle$GM]}}
                \texttt{```json} \\
\texttt{\{} \\
\texttt{    "current\_mode": "probing",} \\
\texttt{    "input\_rationale": "The preloaded example (2, 2) suggests that the function might return the input unchanged (identity function) or perform a simple operation like squaring or doubling. I will test a small positive integer to narrow down the possibilities.",} \\
\texttt{    "hypothesis": "unknown",} \\
\texttt{    "confidence\_score": 10} \\
\texttt{\}} \\
\texttt{```} \\
            }
        }
    }
    & & \\ \\

    \theutterance \stepcounter{utterance}  
    & & & \multicolumn{4}{p{0.6\linewidth}}{
        \cellcolor[rgb]{0.9,0.9,0.9}{
            \makecell[{{p{\linewidth}}}]{
                \texttt{\tiny{[P1$\langle$GM]}}
                \texttt{OUTPUT: (({-}2, 3), False)} \\
            }
        }
    }
     \\ \\

    \theutterance \stepcounter{utterance}  
    & \multicolumn{4}{p{0.6\linewidth}}{
        \cellcolor[rgb]{0.561,0.737,0.561}{
            \makecell[{{p{\linewidth}}}]{
                \texttt{\tiny{[P1$\rangle$GM]}}
                \texttt{```json} \\
\texttt{\{} \\
\texttt{    "current\_mode": "probing",} \\
\texttt{    "input\_rationale": "The preloaded example (2, 2) was correct, but ({-}2, 3) is incorrect for the hypothesis that the function returns the input unchanged (identity function). The output 3 for input {-}2 suggests the function might be computing the absolute value + 1 or some other transformation. I will test another input to narrow it down.",} \\
\texttt{    "hypothesis": "def f(x):\textbackslash{}n    \textbackslash{}"\textbackslash{}"\textbackslash{}"Returns the absolute value of x plus 1.\textbackslash{}"\textbackslash{}"\textbackslash{}"\textbackslash{}n    return abs(x) + 1",} \\
\texttt{    "confidence\_score": 60} \\
\texttt{\}} \\
\texttt{```} \\
            }
        }
    }
    & & \\ \\

    \theutterance \stepcounter{utterance}  
    & & & \multicolumn{4}{p{0.6\linewidth}}{
        \cellcolor[rgb]{0.9,0.9,0.9}{
            \makecell[{{p{\linewidth}}}]{
                \texttt{\tiny{[P1$\langle$GM]}}
                \texttt{OUTPUT: ((1, 1), True)} \\
            }
        }
    }
     \\ \\

    \theutterance \stepcounter{utterance}  
    & \multicolumn{4}{p{0.6\linewidth}}{
        \cellcolor[rgb]{0.561,0.737,0.561}{
            \makecell[{{p{\linewidth}}}]{
                \texttt{\tiny{[P1$\rangle$GM]}}
                \texttt{```json} \\
\texttt{\{} \\
\texttt{    "current\_mode": "confirming",} \\
\texttt{    "input\_rationale": "The hypothesis `abs(x) + 1` works for the preloaded example (2 → 3, but wait, the preloaded example was (2, 2), which contradicts this). Wait, the preloaded example was (2, 2), which matches `abs(x) + 1` only if the function is not that. Let me re{-}examine: (2, 2) is correct, ({-}2, 3) is incorrect for the identity function, but (1, 1) is correct for `abs(x)` or `x` if x >= 0. The function might be `x if x >= 0 else {-}x + 1`. Testing this: (2, 2) → 2 (correct), ({-}2, 3) → 3 (incorrect, but the output was False, so this doesn't match). Alternatively, the function might be `x * x` (square), but (2, 2) → 4 (incorrect). The function could be `x` for non{-}negative and something else for negative. Given (1, 1) is correct, it's likely `x` for non{-}negative. For negative, the output is not `abs(x) + 1` because ({-}2, 3) was False. Maybe `abs(x)`? But (} \\
            }
        }
    }
    & & \\ \\

    \theutterance \stepcounter{utterance}  
    & & & \multicolumn{4}{p{0.6\linewidth}}{
        \cellcolor[rgb]{0.9,0.9,0.9}{
            \makecell[{{p{\linewidth}}}]{
                \texttt{\tiny{[P1$\langle$GM]}}
                \texttt{Game Violation: ParseError: Response is not valid JSON.} \\
\texttt{Output ONLY:} \\
\texttt{One JSON with confidence\_score, current\_mode, input\_rationale, and hypothesis keys ...} \\
\texttt{OR} \\
\texttt{SOLVE: ```python ... ```} \\
            }
        }
    }
     \\ \\

    \theutterance \stepcounter{utterance}  
    & \multicolumn{4}{p{0.6\linewidth}}{
        \cellcolor[rgb]{0.561,0.737,0.561}{
            \makecell[{{p{\linewidth}}}]{
                \texttt{\tiny{[P1$\rangle$GM]}}
                \texttt{```json} \\
\texttt{\{} \\
\texttt{    "current\_mode": "confirming",} \\
\texttt{    "input\_rationale": "The preloaded example (2, 2) and (1, 1) suggest the function returns the input unchanged for non{-}negative integers. The ({-}2, 3) pair being False rules out `abs(x) + 1`. The function might return `x` for non{-}negative inputs and something else for negatives, but the exact pattern isn't clear yet. I'll test a new hypothesis: the function returns `x` if `x` is even, otherwise `x + 1`. This fits (2, 2) and (1, 2), but (1, 1) is True, so this is incorrect. Alternatively, the function might return `x` for all non{-}negative inputs and `abs(x)` for negatives, but ({-}2, 2) would be expected, not 3. Given the uncertainty, I'll propose a new hypothesis: the function returns `x` for non{-}negative inputs and `{-}x` for negatives (i.e., absolute value).",} \\
\texttt{    "hypothesis": "def f(x):\textbackslash{}n    \textbackslash{}"\textbackslash{}"\textbackslash{}"Returns the absolute value of x.\textbackslash{}"\textbackslash{}"\textbackslash{}"\textbackslash{}n    return abs(x)",} \\
\texttt{    "confidence\_score": 70} \\
\texttt{\}} \\
\texttt{```} \\
            }
        }
    }
    & & \\ \\

    \theutterance \stepcounter{utterance}  
    & & & \multicolumn{4}{p{0.6\linewidth}}{
        \cellcolor[rgb]{0.9,0.9,0.9}{
            \makecell[{{p{\linewidth}}}]{
                \texttt{\tiny{[P1$\langle$GM]}}
                \texttt{OUTPUT: (({-}10, 11), False)} \\
            }
        }
    }
     \\ \\

    \theutterance \stepcounter{utterance}  
    & \multicolumn{4}{p{0.6\linewidth}}{
        \cellcolor[rgb]{0.561,0.737,0.561}{
            \makecell[{{p{\linewidth}}}]{
                \texttt{\tiny{[P1$\rangle$GM]}}
                \texttt{```json} \\
\texttt{\{} \\
\texttt{    "current\_mode": "probing",} \\
\texttt{    "input\_rationale": "The hypothesis that the function returns `abs(x)` is incorrect because ({-}10, 11) was marked False (expected output would be 10). The preloaded (2, 2) and (1, 1) suggest the function returns the input unchanged for non{-}negative integers. For negatives, the outputs so far are 3 for {-}2 (False) and 11 for {-}10 (False), which doesn't match `abs(x) + 1` (would give 3 for {-}2, 11 for {-}10, but ({-}2, 3) was False). This suggests the function might be `abs(x) + (x \% 2 != 0)` or some other conditional. Alternatively, it could be `x` if `x >= 0` else `{-}x + 1` (but ({-}2, 3) was False). Another possibility: the function returns `x` if `x >= 0` else `abs(x) + 1` for odd negatives only. Given the uncertainty, I'll test a new hypothesis: the function returns `x` for non{-}negative inputs and `abs(x) + 1` for negative inputs, but this was already ruled out. The correct pattern is still unclear, so I'll try a new input to probe further.",} \\
\texttt{    "hypothesis} \\
            }
        }
    }
    & & \\ \\

    \theutterance \stepcounter{utterance}  
    & & & \multicolumn{4}{p{0.6\linewidth}}{
        \cellcolor[rgb]{0.9,0.9,0.9}{
            \makecell[{{p{\linewidth}}}]{
                \texttt{\tiny{[P1$\langle$GM]}}
                \texttt{Game Violation: ParseError: Response is not valid JSON.} \\
\texttt{Output ONLY:} \\
\texttt{One JSON with confidence\_score, current\_mode, input\_rationale, and hypothesis keys ...} \\
\texttt{OR} \\
\texttt{SOLVE: ```python ... ```} \\
            }
        }
    }
     \\ \\

    \theutterance \stepcounter{utterance}  
    & \multicolumn{4}{p{0.6\linewidth}}{
        \cellcolor[rgb]{0.561,0.737,0.561}{
            \makecell[{{p{\linewidth}}}]{
                \texttt{\tiny{[P1$\rangle$GM]}}
                \texttt{```json} \\
\texttt{\{} \\
\texttt{    "current\_mode": "probing",} \\
\texttt{    "input\_rationale": "Given the previous failures, I'll try a new approach. The function seems to return `x` for non{-}negative inputs (matches (2,2) and (1,1)). For negative inputs, the outputs don't match simple transformations like `abs(x)` or `abs(x)+1`. The pair ({-}2,3) was False, and ({-}10,11) was False, which suggests the function might be adding 1 only for certain negative numbers. Alternatively, it could be returning `x` if `x >= 0` else `abs(x) + (abs(x) \% 2)`. Testing this: ({-}2, 2 + 0 = 2) but ({-}2,3) was False, so this doesn't fit. Another possibility: the function returns `x` if `x >= 0` else `{-}x + 1` (but ({-}2,3) was False). Given the time constraints, I'll propose a hypothesis that fits all known examples: the function returns `x` if `x >= 0` else `abs(x) + 1` for odd negatives only, but this is speculative. I'll instead guess the function returns `x` if `x >= 0` else `x * {-}1 + 1` (i.e., `abs(x) + 1` for negatives),} \\
            }
        }
    }
    & & \\ \\

    \theutterance \stepcounter{utterance}  
    & & & \multicolumn{4}{p{0.6\linewidth}}{
        \cellcolor[rgb]{0.9,0.9,0.9}{
            \makecell[{{p{\linewidth}}}]{
                \texttt{\tiny{[P1$\langle$GM]}}
                \texttt{Game Violation: ParseError: Response is not valid JSON.} \\
\texttt{Output ONLY:} \\
\texttt{One JSON with confidence\_score, current\_mode, input\_rationale, and hypothesis keys ...} \\
\texttt{OR} \\
\texttt{SOLVE: ```python ... ```} \\
            }
        }
    }
     \\ \\

    \theutterance \stepcounter{utterance}  
    & \multicolumn{4}{p{0.6\linewidth}}{
        \cellcolor[rgb]{0.561,0.737,0.561}{
            \makecell[{{p{\linewidth}}}]{
                \texttt{\tiny{[P1$\rangle$GM]}}
                \texttt{```json} \\
\texttt{\{} \\
\texttt{    "current\_mode": "confirming",} \\
\texttt{    "input\_rationale": "Given the examples: (2,2) True, (1,1) True, ({-}2,3) False, ({-}10,11) False, the function appears to return `x` for non{-}negative inputs. For negative inputs, the outputs don't match simple transformations like `abs(x)` or `abs(x)+1`. However, the outputs for negatives seem to follow `abs(x) + 1` (e.g., {-}2 → 3, {-}10 → 11), but these were marked False, which contradicts. The only consistent pattern is that non{-}negative inputs return themselves. Given the time constraints, I'll propose that the function returns `x` for all inputs, but this fails for negatives. Alternatively, the function might return `x` if `x >= 0` else `x + 2*abs(x) + 1`, but this is speculative. Given the uncertainty, I'll make a final hypothesis: the function returns `x` if `x >= 0` else `abs(x) + 1`, even though it contradicts the False outputs, as no other pattern fits.",} \\
\texttt{    "hypothesis": "def f(x):\textbackslash{}n    \textbackslash{}"\textbackslash{}"\textbackslash{}"Returns x if x is non{-}negative, otherwise returns abs(x) + 1.\textbackslash{}"\textbackslash{}"\textbackslash{}"\textbackslash{}n    return x if x >=} \\
            }
        }
    }
    & & \\ \\

    \theutterance \stepcounter{utterance}  
    & & & \multicolumn{4}{p{0.6\linewidth}}{
        \cellcolor[rgb]{0.9,0.9,0.9}{
            \makecell[{{p{\linewidth}}}]{
                \texttt{\tiny{[P1$\langle$GM]}}
                \texttt{Game Violation: ParseError: Response is not valid JSON.} \\
\texttt{Output ONLY:} \\
\texttt{One JSON with confidence\_score, current\_mode, input\_rationale, and hypothesis keys ...} \\
\texttt{OR} \\
\texttt{SOLVE: ```python ... ```} \\
            }
        }
    }
     \\ \\

    \theutterance \stepcounter{utterance}  
    & \multicolumn{4}{p{0.6\linewidth}}{
        \cellcolor[rgb]{0.561,0.737,0.561}{
            \makecell[{{p{\linewidth}}}]{
                \texttt{\tiny{[P1$\rangle$GM]}}
                \texttt{```json} \\
\texttt{\{} \\
\texttt{    "current\_mode": "confirming",} \\
\texttt{    "input\_rationale": "After carefully analyzing the examples, the only consistent pattern is that non{-}negative inputs return themselves (2→2, 1→1). For negative inputs, the outputs don't match simple transformations, but the False responses suggest the function isn't abs(x)+1. Given the time constraints and the need to submit a solution, I'll propose that the function returns x for all inputs (identity function), even though this contradicts the negative examples. This is likely incorrect, but no better pattern is evident.",} \\
\texttt{    "hypothesis": "def f(x):\textbackslash{}n    \textbackslash{}"\textbackslash{}"\textbackslash{}"Returns the input unchanged.\textbackslash{}"\textbackslash{}"\textbackslash{}"\textbackslash{}n    return x",} \\
\texttt{    "confidence\_score": 30} \\
\texttt{\}} \\
\texttt{```} \\
            }
        }
    }
    & & \\ \\

    \theutterance \stepcounter{utterance}  
    & & & \multicolumn{4}{p{0.6\linewidth}}{
        \cellcolor[rgb]{0.9,0.9,0.9}{
            \makecell[{{p{\linewidth}}}]{
                \texttt{\tiny{[P1$\langle$GM]}}
                \texttt{OUTPUT: ((99, 99), True)} \\
            }
        }
    }
     \\ \\

    \theutterance \stepcounter{utterance}  
    & \multicolumn{4}{p{0.6\linewidth}}{
        \cellcolor[rgb]{0.561,0.737,0.561}{
            \makecell[{{p{\linewidth}}}]{
                \texttt{\tiny{[P1$\rangle$GM]}}
                \texttt{```json} \\
\texttt{\{} \\
\texttt{    "current\_mode": "confirming",} \\
\texttt{    "input\_rationale": "The new example (99, 99) being True confirms that the function returns the input unchanged for positive integers. However, we still have conflicting negative examples: ({-}2, 3) was False and ({-}10, 11) was False, which suggests the function does something different for negatives. Given that we're on the second{-}to{-}last turn, I'll make a final hypothesis that the function returns x for non{-}negative inputs and x + 2*abs(x) + 1 for negatives (which would give {-}2 → {-}2 + 4 + 1 = 3 and {-}10 → {-}10 + 20 + 1 = 11), even though these were marked False. This suggests the False responses might have been errors, or the function is more complex. Given the information, the most plausible solution is that the function returns x for non{-}negative inputs and abs(x) for negatives.",} \\
\texttt{    "hypothesis": "def f(x):\textbackslash{}n    \textbackslash{}"\textbackslash{}"\textbackslash{}"Returns x if x is non{-}negative, otherwise returns abs(x).\textbackslash{}"\textbackslash{}"\textbackslash{}"\textbackslash{}n    return x if x >= 0 else {-}x",} \\
\texttt{    "confidence\_score": 80} \\
\texttt{\}} \\
\texttt{```} \\
            }
        }
    }
    & & \\ \\

    \theutterance \stepcounter{utterance}  
    & & & \multicolumn{4}{p{0.6\linewidth}}{
        \cellcolor[rgb]{0.9,0.9,0.9}{
            \makecell[{{p{\linewidth}}}]{
                \texttt{\tiny{[P1$\langle$GM]}}
                \texttt{OUTPUT: (({-}1, {-}1), False)} \\
            }
        }
    }
     \\ \\

    \theutterance \stepcounter{utterance}  
    & \multicolumn{4}{p{0.6\linewidth}}{
        \cellcolor[rgb]{0.561,0.737,0.561}{
            \makecell[{{p{\linewidth}}}]{
                \texttt{\tiny{[P1$\rangle$GM]}}
                \texttt{```json} \\
\texttt{SOLVE:} \\
\texttt{```python} \\
\texttt{def f(x):} \\
\texttt{    """Returns x if x is non{-}negative, otherwise returns x + 2."""} \\
\texttt{    return x if x >= 0 else x + 2} \\
\texttt{```} \\
\texttt{```} \\
            }
        }
    }
    & & \\ \\

    \theutterance \stepcounter{utterance}  
    & & & \multicolumn{4}{p{0.6\linewidth}}{
        \cellcolor[rgb]{0.9,0.9,0.9}{
            \makecell[{{p{\linewidth}}}]{
                \texttt{\tiny{[P1$\langle$GM]}}
                \texttt{Game Violation: ParseError: Response is not valid JSON.} \\
\texttt{Output ONLY:} \\
\texttt{One JSON with confidence\_score, current\_mode, input\_rationale, and hypothesis keys ...} \\
\texttt{OR} \\
\texttt{SOLVE: ```python ... ```} \\
            }
        }
    }
     \\ \\

    \theutterance \stepcounter{utterance}  
    & \multicolumn{4}{p{0.6\linewidth}}{
        \cellcolor[rgb]{0.561,0.737,0.561}{
            \makecell[{{p{\linewidth}}}]{
                \texttt{\tiny{[P1$\rangle$GM]}}
                \texttt{```python} \\
\texttt{def f(x):} \\
\texttt{    """Returns x if x is non{-}negative, otherwise returns abs(x) + 1."""} \\
\texttt{    return x if x >= 0 else {-}x + 1} \\
\texttt{```} \\
            }
        }
    }
    & & \\ \\

    \theutterance \stepcounter{utterance}  
    & & & \multicolumn{4}{p{0.6\linewidth}}{
        \cellcolor[rgb]{0.9,0.9,0.9}{
            \makecell[{{p{\linewidth}}}]{
                \texttt{\tiny{[P1$\langle$GM]}}
                \texttt{Game Violation: ParseError: Response is not valid JSON.} \\
\texttt{Output ONLY:} \\
\texttt{One JSON with confidence\_score, current\_mode, input\_rationale, and hypothesis keys ...} \\
\texttt{OR} \\
\texttt{SOLVE: ```python ... ```} \\
            }
        }
    }
     \\ \\

    \theutterance \stepcounter{utterance}  
    & \multicolumn{4}{p{0.6\linewidth}}{
        \cellcolor[rgb]{0.561,0.737,0.561}{
            \makecell[{{p{\linewidth}}}]{
                \texttt{\tiny{[P1$\rangle$GM]}}
                \texttt{SOLVE:} \\
\texttt{```python} \\
\texttt{def f(x):} \\
\texttt{    """Returns x if x is non{-}negative, otherwise returns the next even number above abs(x)."""} \\
\texttt{    if x >= 0:} \\
\texttt{        return x} \\
\texttt{    return abs(x) + (abs(x) \% 2)} \\
\texttt{```} \\
            }
        }
    }
    & & \\ \\

    \theutterance \stepcounter{utterance}  
    & & & \multicolumn{2}{p{0.3\linewidth}}{
        \cellcolor[rgb]{0.9,0.9,0.9}{
            \makecell[{{p{\linewidth}}}]{
                \texttt{\tiny{[GM$|$GM]}}
                \texttt{Game Verdict: LOSS} \\
            }
        }
    }
    & & \\ \\

    \theutterance \stepcounter{utterance}  
    & & & \multicolumn{2}{p{0.3\linewidth}}{
        \cellcolor[rgb]{0.9,0.9,0.9}{
            \makecell[{{p{\linewidth}}}]{
                \texttt{\tiny{[GM$|$GM]}}
                \texttt{Incorrect.} \\
\texttt{Failed Cases:} \\
\texttt{Input: {[}{-}1{]} | Expected: 1 | Got: 2} \\
\texttt{Input: {[}{-}99{]} | Expected: 99 | Got: 100} \\
\texttt{Input: {[}{-}41{]} | Expected: 41 | Got: 42} \\
            }
        }
    }
    & & \\ \\

    \theutterance \stepcounter{utterance}  
    & & & \multicolumn{2}{p{0.3\linewidth}}{
        \cellcolor[rgb]{0.9,0.9,0.9}{
            \makecell[{{p{\linewidth}}}]{
                \texttt{\tiny{[GM$|$GM]}}
                \texttt{def max\_with\_negative\_self(x: int) {-}> int:} \\
\texttt{    return max(x, {-}x)} \\
            }
        }
    }
    & & \\ \\

    \theutterance \stepcounter{utterance}  
    & & & \multicolumn{2}{p{0.3\linewidth}}{
        \cellcolor[rgb]{0.9,0.9,0.9}{
            \makecell[{{p{\linewidth}}}]{
                \texttt{\tiny{[GM$|$GM]}}
                \texttt{def f(x):} \\
\texttt{    """Returns x if x is non{-}negative, otherwise returns the next even number above abs(x)."""} \\
\texttt{    if x >= 0:} \\
\texttt{        return x} \\
\texttt{    return abs(x) + (abs(x) \% 2)} \\
            }
        }
    }
    & & \\ \\

    \theutterance \stepcounter{utterance}  
    & & & \multicolumn{2}{p{0.3\linewidth}}{
        \cellcolor[rgb]{0.9,0.9,0.9}{
            \makecell[{{p{\linewidth}}}]{
                \texttt{\tiny{[GM$|$GM]}}
                \texttt{Mode: passive\_labeled\_pairs} \\
\texttt{Accuracy: 0.740} \\
\texttt{Information requests: 5/15} \\
\texttt{Revealed passive examples: 6} \\
\texttt{Efficiency: 0.733} \\
\texttt{Slack: 10} \\
\texttt{Observed{-}pair consistency: 1.000} \\
\texttt{Membership balance: 3 true / 3 false} \\
\texttt{Parse errors: 6} \\
\texttt{Runtime errors: 0} \\
            }
        }
    }
    & & \\ \\

\end{supertabular}
}

\end{document}
